\section{A plan for the graphical interface}
\label{app:gui}

This appendix plans the desktop interface that sits on top of the engine.  The
engine was built headless precisely so that this work could be added without
rewriting it, and the plan below leans on what is already there.

\textbf{Status.}  Milestones M0 to M3 are implemented in the \texttt{rubiks\_cfop.gui}
package and start with \texttt{rubiks-cfop-gui} (or \texttt{python -m
rubiks\_cfop.gui}): the scramble bar with random scrambles and seeds, the net
canvas drawn through a painter, the solution table and a move-by-move player.  From
M4, a superseded job's result is ignored rather than the solve being canceled,
because \texttt{solve\_cube} does not take a cancellation token yet.  The plan for a
3D view is in \texttt{docs/gui-3d-plan.md}.

\subsection{Goals and non-goals}

\textbf{Goals.}

\begin{enumerate}[leftmargin=1.4em]
  \item Show the scrambled cube as a clear, colored net, and let the user read a
  solution stage by stage.
  \item Accept a scramble typed by the user, a facelet string from another tool, or
  a generated scramble with a reproducible seed.
  \item Let the user choose the cross color, or let the program choose the best one
  and say what it chose.
  \item Animate a solution move by move, with play, pause, step and a stage scrubber.
  \item Stay responsive: never block the window while solving, and always allow the
  solve to be canceled.
  \item Present errors as information, not as a crash: a bad token should be shown
  at its position in the input.
\end{enumerate}

\textbf{Non-goals.}

\begin{itemize}[leftmargin=1.4em]
  \item A photorealistic 3D cube.  A 2D net is unambiguous, cheap to draw and much
  easier to test; a 3D view is a later, optional renderer behind the same interface.
  \item Competition features: no timer, no inspection countdown, no session
  statistics, no scrambles from the network.
  \item Other puzzles, and any change to the solving algorithm.
\end{itemize}

\subsection{What the engine already provides}

Every one of these exists today and was designed for this purpose:

\begin{longtable}{@{}p{4.8cm}p{9.4cm}@{}}
\toprule
Engine facility & What the interface does with it \\
\midrule
\endhead
\texttt{Cube.apply\_scramble} & turns the typed scramble into a state, raising
\texttt{NotationError} with the token and index to highlight \\
\texttt{solve\_cube} & the whole CFOP solution, already verified, with the cross
color, the rotation applied and every stage \\
\texttt{Solution.to\_dict} & the payload the interface renders: stage labels, move
lists, and the facelets after each stage \\
\texttt{Stage.metadata} & the OLL number, the PLL name, the F2L slot and case, for
informative rows \\
\texttt{cross\_options} & the six crosses with their lengths, for the picker \\
\texttt{net\_layout} & a structured grid of colors; no escape codes to parse, and
no dependency on the terminal \\
\texttt{NetLayout.to\_dict} & the six faces as color names, ready for a web view \\
\texttt{DEFAULT\_PALETTE} & the standard colors, so the canvas matches the CLI \\
\texttt{CancelToken} & canceling a solve from the interface thread \\
\texttt{preload} & building the scrambler tables at start-up, off the UI thread \\
\texttt{RubiksError.to\_dict} & structured errors to display without parsing prose \\
\bottomrule
\end{longtable}

\subsection{Small engine additions the interface needs}

Four additions, all backwards compatible and all testable without a display:

\begin{enumerate}[leftmargin=1.4em]
  \item \texttt{iter\_move\_states(cube, moves)} --- an iterator over
  \texttt{(move, state)} pairs, so the animation applies one move at a time instead
  of owning a copy of the cube model.  Trivial to implement, and it keeps move
  application in one place.
  \item \texttt{Session} --- a stepwise solver: \texttt{next\_stage()} returns the
  next \texttt{Stage}, with \texttt{state}, \texttt{done} and \texttt{progress}
  properties.  A one-shot solve is then just a loop over it, so the interface can
  show a stage as soon as it is ready.
  \item \texttt{on\_stage} callback on the solve options, called from the solving
  thread as each stage completes.  Documented as \emph{called from that thread}:
  the interface posts to its own queue.
  \item \texttt{schemas/solution.schema.json} and a \texttt{schema\_version} field in
  the payload, plus a contract test that validates the emitted JSON against the
  schema.  This is what lets a second front-end (the web one) rely on the format.
\end{enumerate}

None of these adds a dependency to the engine, and none changes an existing
signature.

\subsection{Architecture}

The interface is split in three, and only the last layer knows about a toolkit:

\begin{lstlisting}[style=data]
+-------------------------------------------------------------+
|  app shell (main window, menu, shortcuts, settings dialog)  |
+-------------------------------------------------------------+
|  views        ScrambleBar  CubeCanvas  SolutionTable        |
|               Player       CrossPicker StatusBar            |
+-------------------------------------------------------------+
|  view-model   AppState: cube, solution, selected stage,     |
|               selected move, busy flag, last error          |
|               pure functions, no toolkit, fully unit tested |
+-------------------------------------------------------------+
|  engine       rubiks_cfop  (already headless, thread-safe)  |
+-------------------------------------------------------------+
\end{lstlisting}

Two rules make the plan work:

\begin{enumerate}[leftmargin=1.4em]
  \item \textbf{The view-model never imports a toolkit.}  It holds the state and the
  transitions --- \emph{start a solve}, \emph{select a stage}, \emph{step a move},
  \emph{report an error} --- and is tested with plain pytest.  Everything that is
  easy to get wrong lives here, where it can be tested headlessly.
  \item \textbf{The canvas draws through a painter interface.}  \texttt{CubeCanvas}
  receives a \texttt{Painter} with \texttt{fill\_rect(x, y, w, h, color)} and
  \texttt{draw\_text(...)}.  Production uses the toolkit's painter, tests use a
  recorder, so the layout can be asserted without a screen.
\end{enumerate}

\subsection{Choosing a toolkit}

\begin{longtable}{@{}p{3cm}p{2.4cm}p{4.2cm}p{4.4cm}@{}}
\toprule
Option & Dependency & Strengths & Costs \\
\midrule
\endhead
Tkinter & none (standard library) & matches the engine's zero-dependency stance;
ships with Python; enough for a net, a table and a canvas & dated look; the
\texttt{Canvas} widget is a little low level \\
PySide6 / Qt & a large wheel & polished widgets, good high-DPI support, animation
timers & a heavy dependency and a license to respect; overkill for a net and a
table \\
Web front-end & a server plus a browser & the JSON payload already suits it;
portable; easy to share & a second process, a second language, and no offline
single-file story \\
\bottomrule
\end{longtable}

\textbf{Recommendation:} Tkinter for the first interface.  It keeps the project
installable with no optional dependency, and the view-model plus painter split means
a Qt or web front-end can replace the view layer later without touching the
reasoning.  A web front-end is the natural second step, and the JSON payload and its
schema are the contract for it.

\subsection{Screens and widgets}

\begin{description}[leftmargin=0em,style=nextline]
  \item[Scramble bar] one line of text in WCA notation, with a \emph{Random}
  button, a seed field and a \emph{lock seed} toggle, a style selector
  (random-state, or random-move with a warning that it is not WCA-legal), and a
  \emph{Solve} button.  Invalid notation is marked at the offending token, taken
  from \texttt{NotationError.index}.
  \item[Cube canvas] the unfolded net, drawn from \texttt{net\_layout} in the chosen
  palette.  It shows either the scramble or the state after the selected stage, and
  it is the surface the animation draws on.
  \item[Solution table] one row per stage: label, case (the OLL number, the PLL name,
  the F2L slot), the moves, the length and the running total.  Selecting a row moves
  the canvas to that stage's \texttt{facelets} and the player to that stage.
  \item[Player] play, pause, step forward and back, a speed slider and a stage
  scrubber.  Playing a stage animates its moves one at a time; stepping shows the
  cube after a single move.
  \item[Cross picker] the six colors with their move counts, best first, from
  \texttt{cross\_options}; picking one re-solves with that cross and reports where
  the cube was rotated.
  \item[Status bar] the version, the time spent, the total move count and a
  \emph{verified} badge that comes from the payload, plus a cancel button while a
  solve runs.
  \item[Settings] palette (standard or a high-contrast set), animation speed, font
  size, and whether a generated scramble is remembered between sessions.
\end{description}

\subsection{Threading, responsiveness and cancellation}

\begin{enumerate}[leftmargin=1.4em]
  \item \textbf{Nothing expensive on the UI thread.}  A CFOP solve takes up to a few
  seconds and building the scrambler tables takes about ten.  Both run on worker
  threads.
  \item \textbf{One solve at a time.}  Starting a second solve cancels the first;
  the token makes that prompt and the \texttt{CancelledError} is reported as
  \emph{canceled}, not as a failure.
  \item \textbf{Results cross threads as data.}  The worker hands back
  \texttt{Solution.to\_dict()}; the UI thread rebuilds the view-model from it.  No
  engine object is shared while it is being computed.
  \item \textbf{Table preloading at start-up.}  A splash or a status message during
  \texttt{preload()}, which then makes every later scramble immediate.
  \item \textbf{Cooperative timers for animation.}  The player advances one move per
  tick and stops when the window is not visible; the engine is only asked for the
  next state, never for a whole animation.
\end{enumerate}

\subsection{Testing without a display}

A graphical interface is usually the least tested part of a project.  The
architecture above makes most of it testable anyway:

\begin{longtable}{@{}p{4.6cm}p{9.4cm}@{}}
\toprule
Layer & How it is tested \\
\midrule
\endhead
view-model & plain pytest: transitions, stage selection, move stepping, error
reporting, busy and canceled states \\
canvas layout & a recording painter: the sequence of rectangles and colors must
match \texttt{net\_layout} cell for cell \\
payload contract & the schema file plus a test that validates
\texttt{Solution.to\_dict()} and the CLI's \texttt{--json} against it \\
engine & the existing suite; the interface adds no logic below the view-model \\
interactive smoke test & a scripted session on a headless display where one is
available, skipped otherwise, that types a scramble, solves, steps twice and
asserts the canvas state \\
\bottomrule
\end{longtable}

The rule of thumb: if a behavior cannot be tested without a screen, it belongs in
the view-model or below, not in the widget.

\subsection{Milestones}

Each milestone leaves the program in a usable state and ends with a green gate
(formattingl, linting and the whole test suite).

\begin{longtable}{@{}p{1.5cm}p{5.4cm}p{7.2cm}@{}}
\toprule
Step & Scope & Acceptance \\
\midrule
\endhead
M0 & Skeleton: window, view-model, canvas with the painter interface, static net &
\texttt{rubiks-cfop-gui --scramble "R U R' U'"} shows the net; view-model and canvas
layout tests green; no change to the engine \\
M1 & Scramble bar, solve button, solution table, verified badge & twenty random
scrambles solved through the interface; the payload equals the CLI's
\texttt{--json}; a bad token is reported at its index \\
M2 & Player and animation & play, pause, step and speed animate a solution;
stepping lands on the exact \texttt{facelets} of each stage \\
M3 & Cross picker, seed field, preload splash & choosing a color re-solves with
it; the same seed reproduces the same scramble; no freeze during preload \\
M4 & Cancellation, progress, structured error dialogs & cancel returns within a
fraction of a second and leaves the window usable; every error class renders \\
M5 & Settings, palette, packaging, documentation & a \texttt{[gui]} extra and a
\texttt{rubiks-cfop-gui} entry point; screenshots in the manual; the suite is still
green without a display \\
\bottomrule
\end{longtable}

\subsection{Packaging}

\begin{lstlisting}[style=shell]
# optional dependency group, so the engine stays dependency free
uv add --optional gui pytest-qt   # or nothing at all, for Tkinter

# entry point
rubiks-cfop-gui --scramble "R U R' U'" --cross white
\end{lstlisting}

With Tkinter the group is empty, which is the strongest argument for starting there:
\texttt{pip install rubiks-cfop} gives a working command line, and
\texttt{pip install 'rubiks-cfop[gui]'} gives the window.

\subsection{Risks}

\begin{longtable}{@{}p{5.4cm}p{8.8cm}@{}}
\toprule
Risk & Mitigation \\
\midrule
\endhead
The toolkit choice turns out to be wrong & the view-model and the painter interface
isolate it; a replacement touches only the view layer \\
Slow solves make the interface feel broken & the two F2L improvements in
Section~\ref{sec:trouble} are the real fix; meanwhile the progress callback and the
cancel button keep it honest \\
Animation looks choppy & redrawing 54 cells per frame is negligible; if needed,
animate only the affected layer \\
The engine API drifts while the interface is written & the payload schema and its
contract test are the interface between them \\
Nothing can be tested on the build machine & the canvas painter and the view-model
need no display; the interactive test is skipped when there is no display \\
\bottomrule
\end{longtable}

\subsection{Order of work}

M0 and M1 are deliberately small: they prove the two risky assumptions --- that the
canvas can be tested headlessly, and that a solve can be run off the UI thread
without freezing --- before any polish.  If those hold, M2 to M5 are ordinary
interface work, and the web front-end becomes a matter of writing a second view over
the same payload.
